\documentclass[10pt]{article}
\usepackage[a4paper,margin=0.72in]{geometry}
\usepackage{amsmath,amssymb,amsthm,mathtools}
\usepackage{graphicx}
\usepackage[hidelinks]{hyperref}
\setlength{\parindent}{0pt}
\setlength{\parskip}{0.45em}
\newtheorem{theorem}{Theorem}
\newtheorem{corollary}[theorem]{Corollary}
\newtheorem{example}[theorem]{Example}
\newcommand{\N}{\mathbb N}
\newcommand{\R}{\mathbb R}
\title{An exact residue criterion for completely alternating\\
rational sequences}
\author{Candidate full solution to Question 1.2 of arXiv:2406.13291}
\date{11 August 2026}

\begin{document}
\maketitle

\textbf{Status.} Candidate full solution, likely valid.  The source's general
partial-sum conditions are upgraded to a necessary-and-sufficient condition.

\section*{Source question}

Bhattacharjee and Nailwal \cite{BN} ask which sequences $\{r_k(n)\}_{n\geq0}$
are completely alternating when
\[
 r_k(x)=\frac{p(x)}{\prod_{i=1}^k(x+b_i)},\qquad
 0<b_1<\cdots<b_k,\qquad \deg p\leq k+1.
\]

\begin{center}
\includegraphics[width=0.97\textwidth]{figures/open_problem_crop.png}
\end{center}

The screenshot is Question 1.2 on page 2 of the source PDF.  The source proves
a sufficient condition in terms of partial sums of the residues, but not a
general necessary-and-sufficient statement.

\section*{Proof intuition}

Take one forward difference.  Complete alternation of $r(n)$ is exactly
complete monotonicity of $d_n=r(n+1)-r(n)$.  Partial fractions express $d_n$
as the moments of one explicit finite signed measure on $[0,1]$.  The Hausdorff
moment theorem says that $d$ is completely monotone precisely when it has a
positive representing measure, and polynomial density makes that measure
unique even among finite signed measures.  Positivity can therefore be read
off directly: the linear term supplies a nonnegative atom at $1$, and the
residue exponential polynomial must have the opposite sign on $(0,1)$.

\section*{Exact classification}

Let $p\in\R[x]$ and put $q(x)=\prod_{i=1}^k(x+b_i)$.  The partial-fraction
decomposition has the unique form
\begin{equation}\label{eq:pf}
 \frac{p(x)}{q(x)}=\beta+\alpha x+\sum_{i=1}^k\frac{c_i}{x+b_i},
 \qquad
 c_i=\frac{p(-b_i)}{\prod_{j\ne i}(b_j-b_i)}.
\end{equation}
Here $\alpha=0$ if $\deg p\leq k$, while for $\deg p=k+1$ it is the leading
coefficient of $p$.  Define
\[
                         w(t)=\sum_{i=1}^k c_i t^{b_i}\qquad(0<t<1).
\]

Use the source convention $\nabla=I-E$, where $(Ea)_n=a_{n+1}$.  Thus a
sequence $a$ is completely alternating when $\nabla^m a_n\leq0$ for all
$m\geq1,n\geq0$.

\begin{theorem}[complete answer to Question 1.2]\label{thm:main}
For $r=p/q$ as above, the sequence $\{r(n)\}_{n\geq0}$ is completely
alternating if and only if
\begin{equation}\label{eq:criterion}
                    \alpha\geq0
                    \quad\text{and}\quad
                    w(t)\leq0\quad\text{for every }0<t<1.
\end{equation}
If ``completely alternating'' is additionally required to mean strictly
positive, as in the source's working convention, one adds the single condition
$r(0)>0$.  For nonnegativity, replace it by $r(0)\geq0$.
\end{theorem}

\begin{proof}
Set $d_n=r(n+1)-r(n)=-\nabla r(n)$.  For every $m\geq0$,
\[
                         \nabla^m d_n=-\nabla^{m+1}r(n).
\]
Consequently $r(n)$ is completely alternating if and only if $d=(d_n)$ is
completely monotone.

For $b>0$,
\[
 \frac1{n+b}=\int_0^1t^{n+b-1}\,dt.
\]
Equation~\eqref{eq:pf} therefore gives
\[
 r(n)=\beta+\alpha n+\int_0^1 t^n\frac{w(t)}t\,dt.
\]
The integral is finite because every $b_i>0$.  Taking the forward difference,
\begin{align*}
 d_n
 &=\alpha+\int_0^1t^n(t-1)\frac{w(t)}t\,dt\\
 &=\int_{[0,1]}t^n\,d\sigma(t),
\end{align*}
where
\begin{equation}\label{eq:sigma}
               d\sigma=\alpha\,d\delta_1
                    -\frac{(1-t)w(t)}t\,dt.
\end{equation}
This is a finite signed measure.

If \eqref{eq:criterion} holds, \eqref{eq:sigma} is positive.  Its moments
$d_n$ form a completely monotone sequence by the Hausdorff moment theorem, so
$r(n)$ is completely alternating.

Conversely, suppose $r(n)$ is completely alternating.  Then $d$ is completely
monotone, so the Hausdorff theorem supplies a finite positive Borel measure
$\mu$ on $[0,1]$ with $d_n=\int t^n\,d\mu(t)$ for every $n\geq0$.  The finite
signed measures $\mu$ and $\sigma$ have identical moments.  Hence they agree
on all polynomials, and by Weierstrass approximation they agree on every
continuous function.  Thus $\mu=\sigma$.  Positivity of the atom at $1$ gives
$\alpha\geq0$, and positivity of the absolutely continuous part gives
$w(t)\leq0$ almost everywhere on $(0,1)$.  Since $w$ is continuous there, the
inequality holds everywhere.

Finally, the first-order condition makes $r(n)$ increasing.  Hence all terms
are strictly positive (respectively nonnegative) exactly when the first one is.
\end{proof}

\section*{Relation to the source condition and an effective test}

Let $C_\ell=\sum_{i=1}^\ell c_i$.  Abel summation gives
\[
 w(t)=\sum_{\ell=1}^{k-1}C_\ell
       \bigl(t^{b_\ell}-t^{b_{\ell+1}}\bigr)+C_k t^{b_k}.
\]
Thus the source hypotheses $C_\ell\leq0$ imply $w(t)\leq0$, recovering its
general sufficient theorem.  They are not necessary.

\begin{example}[strict extension of the source test]\label{ex:strict}
Take $b=(1,2,3)$ and $c=(-1,2,-1)$.  Then
\[
 w(t)=-t+2t^2-t^3=-t(1-t)^2\leq0,
\]
although $C_2=1>0$.  These residues arise from
\[
 r(x)=1-\frac{2}{(x+1)(x+2)(x+3)}.
\]
Since $r(0)=2/3>0$, Theorem~\ref{thm:main} gives a positive completely
alternating sequence that is not detected by the source partial-sum criterion.
Indeed,
\[
 r(n+1)-r(n)=\int_0^1t^n(1-t)^3\,dt.
\]
\end{example}

\begin{corollary}[exact finite procedure for rational data]
If $b_i\in\mathbb Q_{>0}$ and $p$ has rational coefficients, the criterion in
Theorem~\ref{thm:main} is decidable by exact rational arithmetic and a Sturm
sequence.
\end{corollary}

\begin{proof}
Choose $L\in\N$ so that $m_i=Lb_i\in\N$.  Substituting $t=u^L$ gives
\[
                         w(t)=P(u):=\sum_i c_i u^{m_i}.
\]
All $c_i$ are rational.  Hence $w\leq0$ on $(0,1)$ exactly when the rational
polynomial $P$ is nonpositive on $(0,1)$, which is decidable by isolating its
real roots (or the roots of $P'$) with Sturm's theorem and checking the finitely
many resulting intervals.
\end{proof}

\section*{Verification, scope, and novelty}

The proof separately checks the atom at $1$, integrability at $0$, and the
source's positivity convention.  Example~\ref{ex:strict} verifies that the
criterion is a genuine strengthening rather than a restatement of the source
partial-sum hypothesis.  No numerical computation is used.

The answer is exact in terms of the residues and poles.  For arbitrary real
exponents it leaves a sign test for a finite exponential polynomial; this is
still a necessary-and-sufficient classification, and for rational data the
corollary makes it a finite exact algorithm.

A bounded novelty search on 11 August 2026 covered the run indexes and
exact-phrase/residue-keyword web searches.  It found the source and published
version, which state only sufficient general conditions, but no later exact
answer.  Novelty confidence is moderate: the mechanism is a short direct
completion of the source's signed-measure calculation and may be implicit
folklore, but the theorem was not located or stated in the source.

\textbf{Human-review recommendation.} Verify the signed Hausdorff-measure
uniqueness step and confirm that Question 1.2 accepts an exact
exponential-polynomial sign criterion rather than requiring a finite list of
ordered zero/pole inequalities.

\begin{thebibliography}{9}
\bibitem{BN}
M. Bhattacharjee and R. Nailwal,
\emph{A characterization of completely alternating functions},
J. Math. Anal. Appl. (2025); arXiv:2406.13291.

\bibitem{BCR}
C. Berg, J. P. R. Christensen, and P. Ressel,
\emph{Harmonic Analysis on Semigroups: Theory of Positive Definite and Related
Functions}, Graduate Texts in Mathematics 100, Springer, 1984.
\end{thebibliography}

\end{document}

